\section{Background and Taxonomy}

Let $g_t=\grad f(x_t)$ and let $C_t$ denote a low-rank projection or
compression operator. For an orthogonal projector $C_t(v)=P_tP_t^\top v$, the
discarded component is $(I-C_t)g_t$. Three mechanisms should be distinguished.

\paragraph{Pure projection.}
Projected optimization updates only $C_t(g_t)$ and discards
$(I-C_t)g_t$. If the subspace is stale, the orthogonal gradient can persist
even while the projected part becomes small.

\paragraph{\Fira-style norm recovery.}
The current gradient is decomposed into a projected branch and a discarded
branch,
\[
R_t=P_t^\top g_t,\qquad g_t^\perp=(I-P_tP_t^\top)g_t.
\]
A low-rank adaptive map produces $\psi_t(R_t)$, and the discarded branch is
reintroduced through a scalar or column-wise factor,
\[
u_t^\parallel = P_t \psi_t(R_t),\qquad
u_t^\perp = s_t g_t^\perp,\qquad
s_t \approx \phi_t(R_t).
\]
This mechanism uses the \emph{current} projected adaptive statistics to scale
the \emph{current} discarded gradient. It does not, by itself, maintain a
residual buffer that guarantees later correction of under-transmitted mass.

\paragraph{Error Feedback.}
Classical Error Feedback forms an accumulator
\[
a_t = g_t + e_t,\qquad c_t = C_t(a_t),\qquad e_{t+1}=a_t-c_t,
\]
and applies the compressed direction $c_t$ to the parameters. The defining
feature is the residual state $e_t$, not merely the presence of a discarded
component in the current step.

\paragraph{\BECR notation and units.}
Throughout the paper, $E_t$, $H_t$, $Z_t$, and $Z_t^{\mathrm{eff}}$ are stored
in raw-gradient units. The scalar $\rho_t$ is a raw-gradient residual
transmission factor, while $s_t$ is an adaptive recovery or preconditioning
scale. The preconditioned orthogonal update is
\[
U_t^\perp = s_t Z_t^{\mathrm{eff}}.
\]
Residual updates never contain the learning rate; the learning rate appears
only in the parameter update.
